
\begin{abstract}
Mixture-of-Experts (MoE) has become the dominant architecture for
frontier language models.
To meet this demand, production frameworks have built optimized
MoE training stacks over years of engineering effort.
Yet evolving these stacks for new architectures and system
optimizations remains expensive.
With the rise of AI coding agents, they could automate parts of
training-framework development and accelerate this evolution.
But applying them to these existing frameworks carries hidden
costs, invisible to today's throughput-only evaluations.
We name this missing dimension agent-task efficiency (ATE): the
cost of using coding agents to understand, operate, and extend a
framework.
Grounded in four agent-native design principles, we build \sys,
a compact, agent-native MoE training framework.
We further introduce \sysbench, covering real-world
training-framework tasks.
Our evaluation shows \sys matches the throughput of production
frameworks, and on \sysbench, \sys enables higher agent-task
efficiency, with up to 62\% fewer Agent Turns and 64\%
less Active GPU Time.

\vspace{1ex}
\noindent \textbf{GitHub repo:} \url{https://github.com/mlc-ai/pith-train}
\end{abstract}
